\begin{abstract}
\noindent
Organisations increasingly expose models rather than data, so a small language model can
often learn only by querying other models, each with its own price and competence. This
project studies provider-selective distillation: given a query budget and several
black-box providers that return text only, how should a student decide whom to ask, about
what, and when? A 1.5B-parameter student is trained on its own responses, rewarded by
agreement with the chosen provider's answers, across three domains with checkable answers,
five providers of different price and quality, and a deliberately degraded copy of the
strongest. A preliminary experiment showed that logit-based distillation across tokenizers
is lossy and unstable, which supports the text-only design. Three hypotheses are tested:
the best provider differs by domain (H1), changes during training (H2), and an adaptive
selector beats the best fixed allocation (H3). So far, cheap providers give the most
learning per unit of budget, and the most expensive one leads only where no provider helps. The proposed selector
matches the best single provider (test utility 0.500 against 0.494) but does not yet beat
it, and differences among the affordable providers are too small to settle H1 and H2. A plan to
settle them follows.

\vspace{5mm}
Subject Descriptors:
\begin{list}{\quad}{}
    \item I.2.6 Learning
    \item I.2.7 Natural Language Processing
    \item I.2.11 Distributed Artificial Intelligence
    \item G.3 Probability and Statistics
\end{list}

Keywords: \\
\tab On-policy distillation, black-box language models, multi-armed bandits, budget
allocation, federated learning

Implementation Software and Hardware:
\begin{list}{\quad}{}
    \item Python 3.12, PyTorch, vLLM, TRL, PEFT, Math-Verify, bubblewrap
    \item NVIDIA L40 GPUs (shared server); SoC Slurm cluster (A100, H100)
\end{list}
\end{abstract}
\pagebreak
